\section{Reading this blueprint}
This document maps the completed QE package to selected mathematical steps
and exact Lean declarations. It is a navigation aid and proof outline;
the linked theorem types and compiler audits determine the formal scope.
Its source links are pinned to the checked proof revision
\texttt{939346b}, verified with Lean 4.35.0 and the pinned Mathlib revision.
\href{https://github.com/daviddhilbert/ACVF-QE/actions/runs/42689540057}{The matching GitHub verification passed}.

The package has three developments: pure ACF, ordinary valued-field ACVF,
and the corrected two-sorted RV language. Trivial and nontrivial valuation
cases are separate. No theorem of uniform QE for their union is claimed.
RV inversion is named, with $0^{-1}=0$; VF inversion remains unnamed.
The RV equality axiom enforces kernel $1+\mathfrak m$, and RV has one
absorbing element. Mathlib's multiplicative valuation order reverses the
paper's value-group order. See the
\repoFile{LANGUAGE-ERRATA.md}{language and theory errata}.

C-minimality is outside this package's claims. Historical interfaces named
after Marker or Engler--Prestel are intermediate propositions: the public
QE proof supplies their relevant inputs by checked proofs. In particular,
\texttt{MarkerProposition4328} is not an extra hypothesis of Theorem 2.3.
The exported QE audits use only \texttt{propext}, \texttt{Classical.choice},
and \texttt{Quot.sound}.


\section{Languages and precise scope}

\begin{definition}[Uniform quantifier elimination]

\label{qe-syntax}

\lean{QECriterion.HasQEb, SpecialTransformations.Milestone0.HasQE}

\leanok

For every formula with parameter indices $\alpha$ and $n$ additional variables,
there is one quantifier-free formula equivalent in every model and at every
assignment. The RV target additionally requires both formulas and assignments
to respect the VF/RV sorts. Parameter sets need not be finite or countable.

\end{definition}

\section{ACF and ordinary valued fields}

\begin{theorem}[Johnson criterion]

\label{johnson}

\lean{HQEDerive.hasQE_of_johnson, HQEDerive.H_of_qfDetermined, BoundedQE.qeExStep_theoryIff, QECriterion.hasQEb_of_QEAllStep}

\leanok

Let $T_0\subseteq T$ in a language of universe level zero. Suppose there are
prime $T$-models over $T_0$-substructures, amalgamation for $T_0$, and
one-variable existential closure of $T$-models in $T_0$-extensions.
Then $T$ admits quantifier elimination.

\end{theorem}

\begin{proof}

\uses{qe-syntax}

\leanok

Quantifier-free types determine one-variable existential formulas.
The checked compactness argument produces uniform formula witnesses.

\end{proof}

\begin{lemma}[Field prime models and amalgamation]

\label{acf-inputs}

\lean{JohnsonACFPrime.hasPrimeModels_field_acf, ACFValidation.amalgamation_field, ACFValidation.oneEC_field_acf}

\leanok

Fields have prime algebraically closed extensions over ring substructures,
field embeddings can be amalgamated, and an algebraically closed field is
one-variable existentially closed in field extensions. No characteristic is fixed.

\end{lemma}

\begin{proof}

\leanok

The prime-model construction and polynomial arguments are supplied by the
checked field-development modules.

\end{proof}

\begin{theorem}[Pure ACF quantifier elimination]

\label{acf}

\lean{JohnsonACF.hasQEb}

\leanok

The ring-language theory of algebraically closed fields admits bounded-formula
quantifier elimination in every characteristic, retaining arbitrary parameter
indices and finite variable contexts.

\end{theorem}

\begin{proof}

\uses{johnson, acf-inputs}

\leanok

Apply the bounded version of the Johnson argument to the three proved
field inputs.

\end{proof}

\begin{lemma}[Valued-field extension inputs]

\label{valued-inputs}

\lean{ValPM.hasPrimeModels_ACVF_val, ValOneEC.hasAmalgamation_valField, ValOneEC.oneECIn_of_modelType}

\leanok

For the ring language with valuation comparison, the checked prime-model,
amalgamation, and one-variable existential-closure constructions supply the
Johnson inputs for nontrivially valued ACVF. No characteristic restriction is imposed.

\end{lemma}

\begin{proof}

\leanok

The proof develops valuation-compatible field maps, residue-root arguments,
and amalgamation in the ordinary valued-field language.

\end{proof}

\begin{theorem}[Nontrivial ordinary ACVF QE]

\label{ordinary}

\lean{hasQE_ACVF_val, hasQEb_ACVF_val}

\leanok

Nontrivially valued algebraically closed fields admit quantifier elimination
in the ring language with the binary valuation-comparison relation.
Both ordinary and bounded-formula interfaces are exported.

\end{theorem}

\begin{proof}

\uses{johnson, valued-inputs}

\leanok

Apply the generic Johnson criterion to the proved valued-field inputs.

\end{proof}

\begin{theorem}[Trivial ordinary ACVF QE]

\label{ordinary-trivial}

\lean{ValPM.hasQE_ACVF_trivial_val, ValPM.hasQEb_ACVF_trivial_val}

\leanok

The separate theory of trivially valued algebraically closed fields admits
ordinary and bounded-formula quantifier elimination in the same valued-field
language.

\end{theorem}

\begin{proof}

\uses{acf}

\leanok

Trivial valuation comparison reduces to zero tests; translate to the
ring language and use pure ACF quantifier elimination.

\end{proof}

\section{Saturation and the generic embedding criterion}

\begin{definition}[Genuine small-parameter saturation]

\label{saturation}

\lean{MarkerBridge.SatAt}

\leanok

A structure $N$ is $\kappa$-saturated when each finitely satisfiable
one-variable partial type over any parameter map $A\to N$, with $|A|<\kappa$,
has a realization in $N$. This definition contains no embedding-extension
or quantifier-elimination conclusion.

\end{definition}

\begin{theorem}[Saturated elementary extension]

\label{ultrapower}

\lean{CountableSaturatedExtension.exists_saturated_elementary_extension}

\leanok

Every nonempty structure in a countable language has a genuinely
$\aleph_1$-saturated elementary extension. Its cardinality need not be
$\aleph_1$, and no continuum hypothesis is assumed.

\end{theorem}

\begin{proof}

\uses{saturation}

\leanok

A countable nonprincipal ultrapower and the checked realization argument
supply saturation. This adapts Elliot Glazer's licensed construction; source
credit and Apache-2.0 notices are retained in the repository.

\end{proof}

\begin{theorem}[Countable-language QE criterion]

\label{sat-criterion}

\lean{CountableQECriterion.fullQE_of_saturatedExtension}

\leanok

For a countable language with a constant, suppose every embedding of a
substructure of a model $M\models T$ into a sufficiently saturated
$N\models T$ extends to an embedding of all of $M$. Then $T$ has full
quantifier elimination. The final models and parameter sets are unrestricted.

\end{theorem}

\begin{proof}

\uses{ultrapower, qe-syntax}

\leanok

Downward L\"owenheim--Skolem reduces the transfer problem to countable
models; the proved ultrapower construction and compactness produce uniform
quantifier-free formulas.

\end{proof}

\begin{theorem}[Embedding extension from QE]

\label{qe-extension}

\lean{QESaturatedExtension.extension_of_fullQE}

\leanok

If $T$ has full quantifier elimination, embeddings from model
substructures extend to the whole source model when the target realizes
types over parameter sets smaller than the successor of the source cardinal.

\end{theorem}

\begin{proof}

\uses{saturation, qe-syntax}

\leanok

Quantifier elimination makes the pushed atomic data finitely satisfiable;
genuine saturation and the checked maximal partial-embedding argument finish
one-step extensions and then the full embedding.

\end{proof}

\section{Preparing and extending RV embeddings}

\begin{theorem}[Residue-field embedding extension]

\label{acf-extension}

\lean{ProvedACFExtension.markerACFSaturatedEmbeddingExtension}

\leanok

Embeddings of subfields of algebraically closed fields extend into genuinely
saturated algebraically closed targets, with the exact ring-language
interpretations used by the RV proof.

\end{theorem}

\begin{proof}

\uses{acf, qe-extension}

\leanok

Convert the proved bounded ACF QE to full QE and apply the checked generic
embedding-extension argument. The former Marker ACF input is discharged.

\end{proof}

\begin{theorem}[Residue fields are algebraically closed]

\label{residue-closure}

\lean{ProvedResidueClosure.provedResidueClosure}

\leanok

For every valuation on an algebraically closed field, its residue field is
algebraically closed. Nontriviality of the valuation is unnecessary.

\end{theorem}

\begin{proof}

\leanok

Reuse the checked residue-root argument from the ordinary valued-field
development. This is a proved theorem, not an additional Engler--Prestel hypothesis.

\end{proof}

\begin{theorem}[Whole valued-field embedding extension]

\label{valued-extension}

\lean{JohnsonValuedExtension.markerNotesProposition45}

\leanok

The valued-subfield extension statement required by the RV argument follows
from ordinary nontrivial ACVF QE and genuine saturation, through the
repository's explicit language-translation bridges.

\end{theorem}

\begin{proof}

\uses{ordinary, qe-extension}

\leanok

Use the proved ordinary valued-field QE, then translate saturation and the
resulting embeddings back to the concrete valued-field interface. The former
Marker-notes premise is discharged.

\end{proof}

\begin{theorem}[Compatible prepared RV frame]

\label{frame}

\lean{SubstructureRVPreparation.exists_frame}

\leanok

Given the actual VF/RV language substructure, a compatible valuation
presentation, a substructure embedding, and a sufficiently saturated target,
construct a residue embedding and a multiplicative, inverse-closed RV graph
extending the old RV map. The residue-closure and ACF-extension inputs are
supplied by the preceding proved results.

\end{theorem}

\begin{proof}

\uses{acf-extension, residue-closure, saturation}

\leanok

Extend the residue-field map and assemble the compatible RV-group data.
Named RV inversion supplies the required substructure closure.

\end{proof}

\begin{theorem}[A field map representing the RV graph]

\label{maximal}

\lean{MaximalRVExtension.exists_field_representing_graph}

\leanok

Given algebraically closed valued fields, a residue map, a functional
compatible RV frame, and a compatible map on a starting subfield, there is
an enlarged subfield map representing every prescribed RV pair.

\end{theorem}

\begin{proof}

\leanok

The residue-algebraic, residue-transcendental, value-algebraic,
value-transcendental, and immediate extension cases feed a maximality
argument. Its algebraic and frame hypotheses remain explicit in this
intermediate theorem.

\end{proof}

\begin{theorem}[All RV preparation inputs discharged]

\label{preparation}

\lean{LocalAssembly.preparation}

\leanok

The intermediate preparation interface is inhabited by a checked
construction from the actual language substructure and genuine saturation.
No independent preparation assumption is left to the public QE caller.

\end{theorem}

\begin{proof}

\uses{frame, acf-extension, residue-closure}

\leanok

Combine the prepared RV frame with the proved residue-closure and
ACF-extension results.

\end{proof}

\begin{theorem}[Canonical model embedding extension]

\label{rv-extension}

\lean{LocalAssembly.embedding_of_preparedFrame, LocalAssembly.canonicalExtension_of_preparation}

\leanok

For canonical nontrivial ACVF models, an embedding of an actual
language substructure extends into a sufficiently saturated target.
The intermediate conditional preparation input is supplied by the preceding
checked preparation theorem.

\end{theorem}

\begin{proof}

\uses{preparation, maximal, valued-extension}

\leanok

Represent the RV graph in a subfield, extend its valued-field embedding,
and use the canonical-structure bridges to recover a full tagged-language map.

\end{proof}

\begin{theorem}[Return to the well-sorted fragment]

\label{sorting}

\lean{CanonicalStructure.hasQE_of_raw_canonical}

\leanok

Raw quantifier elimination on the canonical tagged presentations yields
quantifier elimination for the intended well-sorted VF/RV fragment in every
model of the nontrivial ACVF theory.

\end{theorem}

\begin{proof}

\uses{qe-syntax}

\leanok

Canonicalize the tagged presentation and sort the resulting quantifier-free
formula. This keeps the required VF/RV parameter and variable sorts.

\end{proof}

\section{RV quantifier-elimination theorems}

\begin{theorem}[Theorem 2.3: nontrivial RV QE]

\label{rv}

\lean{LocalAssembly.fullQE, SpecialTransformations.Milestone0.theorem_2_3_quantifier_elimination}

\leanok

The nontrivially valued theory in the corrected RV language admits
quantifier elimination on the well-sorted fragment. The output is well-sorted
and uniform over all models and assignments, with arbitrary parameter
indices. The exported theorem has no additional mathematical hypotheses.

\end{theorem}

\begin{proof}

\uses{sat-criterion, rv-extension, sorting}

\leanok

Apply the proved countable-language criterion to the constructed canonical
embedding extension, then pass through the canonical and sorting bridges.

\end{proof}

\begin{theorem}[Canonical trivial RV QE]

\label{rv-trivial}

\lean{SpecialTransformations.TrivialRV.canonical_trivial_rv_qe}

\leanok

For a canonical trivially valued RV structure $K\sqcup K$ of any
algebraically closed field $K$, every tagged formula has a well-sorted,
quantifier-free equivalent, uniform over these structures in every
characteristic and over sorted parameter and variable assignments.

\end{theorem}

\begin{proof}

\uses{acf, qe-syntax}

\leanok

Encode the tags by ring-language coordinates, apply ACF QE, and reify
the quantifier-free formula back into the sorted RV language.

\end{proof}

\begin{theorem}[QE for the common trivial RV theory]

\label{rv-trivial-theory}

\lean{SpecialTransformations.TrivialRV.hasQE}

\leanok

The common sentence theory of canonical trivially valued algebraically
closed RV structures admits the well-sorted quantifier-elimination property.
This is a separate theory from nontrivial ACVF.

\end{theorem}

\begin{proof}

\uses{rv-trivial}

\leanok

Universally close the sorted equivalence on canonical models. The resulting
sentence belongs to their common theory and therefore transfers to every
model of that theory.

\end{proof}
